% abstract
%
\bgroup
\color{abs}
\hrule
\egroup



%
%%  DO NOT EDIT  LINES ABOVE                     %

\section{Resumen}

El presente documento expone el proyecto final del curso IE-0523 Sistemas Digitales II, cuyo objetivo principal consiste en la implementación en Verilog HDL de la subcapa de codificación física (\textit{Physical Coding Sublayer}, PCS) definida en la cláusula 36 del estándar IEEE 802.3 para el tipo 1000BASE-X.

El trabajo desarrollado aborda el análisis detallado del estándar, la descomposición funcional del PCS y el diseño modular de los bloques que conforman la etapa de transmisión (\texttt{TRANSMIT}), sincronización (\texttt{SYNCHRONIZATION}) y recepción (\texttt{RECEIVE}).
Se describen las estructuras arquitectónicas diseñadas, el comportamiento esperado de cada módulo, las señales relevantes y la metodología utilizada para el desarrollo de las pruebas y en el trabajo colaborativo.
Asimismo, se conectaron los tres módulos en modo \textit{loopback} para su prueba final.

El proyecto se basa en un esquema modular, donde se reutilizan componentes y convenciones comunes con el fin de asegurar cohesión y claridad en el desarrollo del código, tanto para el entendimiento del profesor al calificar como para los estudiantes durante la realización del proyecto.
El resultado más importante es que se tuvieron los resultados esperados según la cláusula del estándar dado, lo cual equivale a un proyecto exitoso.
